\subsubsection{Laiko eilučių tinklas (1D CNN)}
\label{subsubsec:keystroke_cnn}

Klasė \texttt{KeystrokeCNN} naudoja \texttt{nn.Conv1d}, \texttt{nn.BatchNorm1d}, \texttt{nn.MaxPool1d} ir \texttt{nn.Dropout} sluoksnius. Įėjimas – $1 \times 31$ formato sekos tenzorius ($31$ klaviatūros laikinių požymių, vienas kanalas). Bazinė konfigūracija turi du konvoliucinius blokus su filtrų skaičiais $(64,\,128)$, branduolio dydžiu $3$ ir telkimo dydžiu $2$. Po dviejų telkimo sluoksnių sekos ilgis sumažėja iki $\lfloor 31 / 2^{2} \rfloor = 7$, todėl po suplojimo gaunamas $128 \cdot 7 = 896$ dydžio vektorius, paduodamas į $128$ neuronų paslėptąjį sluoksnį, o galiausiai – į $30$ neuronų išėjimo sluoksnį (po vieną neuroną kiekvienam dalyviui). Sluoksnių parametrai pateikti \ref{tab:keystroke_cnn_architecture} lentelėje.

\begin{table}[H]
    \centering
    \caption{Laiko eilučių klasifikavimo tinklo (\texttt{KeystrokeCNN}) bazinės architektūros sluoksniai.}
    \begin{tabular}{|c|l|l|}
        \hline
        Nr. & Sluoksnis & Parametrai / išėjimas \\
        \hline
        $0$ & Įėjimas & $1 \times 31$ \\ \hline
        $1$ & \texttt{Conv1d} & $1 \to 64$, branduolys $3$, \texttt{padding}~$=1$ \\ \hline
        $2$ & \texttt{ReLU} & – \\ \hline
        $3$ & \texttt{MaxPool1d} & dydis $2$, išėjimas $64 \times 15$ \\ \hline
        $4$ & \texttt{Conv1d} & $64 \to 128$, branduolys $3$, \texttt{padding}~$=1$ \\ \hline
        $5$ & \texttt{ReLU} & – \\ \hline
        $6$ & \texttt{MaxPool1d} & dydis $2$, išėjimas $128 \times 7$ \\ \hline
        $7$ & \texttt{Flatten} & $896$ \\ \hline
        $8$ & \texttt{Linear} & $896 \to 128$ \\ \hline
        $9$ & \texttt{ReLU} & – \\ \hline
        $10$ & \texttt{Dropout} & $p = 0{,}0$ (bazinė) \\ \hline
        $11$ & \texttt{Linear} & $128 \to 30$ \\ \hline
    \end{tabular}
    \label{tab:keystroke_cnn_architecture}
\end{table}

